% 调用上级目录公共模板
\documentclass{../template/softeng}

\begin{document}
	
	% 封面：文档名、项目名、版本、编写人、审核人、审批人、日期
	\doccover{概要设计说明书}
	{在线OJ在线判题系统}
	{V2026.07}
	
	% 目录
	\tableofcontents
	\newpage
	
	% 修订记录
	\begin{revrecord}
		V1.0 & 2026-07-11 & 张三 & 李四 & 初稿完成，完成系统总体架构、模块、接口、数据结构概要设计梳理 \\
		\hline
		V1.1 & 2026-07-12 & 张三 & 王五 & 补充判题流程、异常出错处理、安全与维护相关设计内容 \\
		\hline
	\end{revrecord}
	
	\section{引言}
	\subsection{编写目的}
	本文档为在线OJ在线判题系统概要设计说明书，基于前期软件需求规格说明书完成系统顶层架构设计。文档用于明确系统整体模块划分、业务处理流程、内外接口规范、数据库数据结构、运行控制方案、异常处理、安全及维护设计，为后续详细设计、编码开发、单元测试提供设计依据。
	本文档阅读对象包含项目负责人、系统架构设计师、前后端开发工程师、测试工程师。
	
	\subsection{项目背景}
	本系统为高校大二软件工程实训开发项目，开发主体为校内实训开发小组，服务于高校程序设计课程教师与在校学生。系统独立部署，无第三方业务系统对接，依靠Web服务、判题沙箱服务、关系型数据库实现完整在线编程、自动评测、教学管理业务闭环，用于支撑线上编程练习、作业自动批改、学情数据统计。
	
	\subsection{定义}
	\begin{itemize}
		\item OJ：Online Judge，在线程序自动判题系统；
		\item AC：Accepted，代码通过全部测试用例；
		\item 判题沙箱：隔离执行用户代码的安全运行环境，限制资源访问，防范恶意代码破坏服务；
		\item 测试用例：包含标准输入、预期输出、边界条件的程序评测数据集；
		\item 判题任务队列：缓存待执行判题任务的异步调度队列；
		\item RESTful API：前后端交互标准接口规范；
	\end{itemize}
	
	\subsection{参考资料}
	\begin{enumerate}
		\item GB/T 9386-2008 计算机软件概要设计说明书规范；
		\item 《在线OJ在线判题系统软件需求规格说明书》V2026.07；
		\item 2026软件工程实训项目任务书；
		\item Web前后端分离开发设计规范；
		\item Linux沙箱安全隔离开发技术文档；
		\item MySQL数据库设计开发规范。
	\end{enumerate}
	
	\section{任务概述}
	\subsection{目标}
	搭建一套前后端分离架构的在线编程自动判题平台，划分学生、教师两类角色；拆解系统为用户模块、题库模块、在线代码编辑与判题模块、提交记录统计模块、讨论区模块、系统监控管理模块；实现题库分级管理、在线代码提交、沙箱自动编译评测、刷题数据排行、师生交流讨论、教师后台运维监控全套能力，满足高校编程课程线上练习、自动批改、学情统计教学需求。
	
	\subsection{运行环境}
	\begin{itemize}
		\item 服务端操作系统：Linux CentOS 7 / Windows Server 2019；
		\item Web服务容器：Nginx、Tomcat；
		\item 数据库：MySQL 8.0；
		\item 判题依赖：GCC、Python、Java等多语言编译器、进程隔离沙箱程序；
		\item 客户端环境：Chrome、Edge、Firefox主流现代浏览器；
		\item 中间件：Redis（缓存、任务队列）。
	\end{itemize}
	
	\subsection{需求概述}
	系统分为学生端与教师管理端两大业务域：
	1. 学生端：账号注册登录、分级题库浏览、在线代码编写提交、自动判题结果反馈、个人提交记录与刷题统计、全站排行榜、题目讨论区发帖交流；
	2. 教师端：题库增删改查、批量导入题目、测试用例配置、全量用户与提交行为监管、异常抄袭提交标记、平台运行状态可视化监控、公告发布、讨论区内容审核管理、教学统计报表导出。
	
	\subsection{条件与限制}
	\begin{enumerate}
		\item 判题模块必须采用沙箱隔离机制，禁止用户代码直接读写服务器文件、访问网络；
		\item 单份代码运行设置时间、内存双重上限，防止服务器资源耗尽；
		\item 系统仅开放学生、教师两类角色，无游客完整答题权限；
		\item 题库批量导入仅支持标准结构化文本文件，不支持自定义复杂格式；
		\item 服务器硬件资源有限，并发判题任务存在最大并发阈值；
		\item 所有接口需携带身份鉴权凭证，严格区分角色数据访问权限。
	\end{enumerate}
	
	\section{总体设计}
	\subsection{处理流程}
	1. 用户注册/登录鉴权，校验身份后跳转对应功能页面；
	2. 学生浏览题库，选定题目进入在线代码编辑器，编写代码提交至后端；
	3. 后端接收代码、题目ID，封装判题任务送入Redis任务队列；
	4. 判题服务消费队列任务，调度沙箱隔离运行代码，加载对应测试用例逐组比对输出；
	5. 收集判题状态、运行耗时、内存占用结果，写入数据库提交记录；
	6. 前端轮询获取判题结果，展示评测状态、错误日志；
	7. 系统同步更新个人解题统计、全站排行榜数据；
	8. 教师登录后台，完成题库维护、用户监管、平台监控、讨论区管理操作；
	9. 所有操作行为记录系统日志，异常数据触发预警提示。
	
	\subsection{总体结构和模块外部设计}
	系统整体采用分层架构，自上而下分为四层：前端展示层、应用服务层、判题服务层、数据持久层，拆分六大独立业务模块：
	\begin{enumerate}
		\item 用户权限模块：账号注册登录、身份鉴权、角色权限控制、个人信息管理；
		\item 题库管理模块：题目CRUD、难度标签管理、测试用例配置、批量导入导出；
		\item 在线判题模块：在线代码编辑器、任务队列调度、沙箱编译运行、结果判定；
		\item 数据统计排行模块：个人提交记录、解题统计、全站实时排行榜、教学报表；
		\item 讨论与公告模块：题目讨论帖子、评论回复、全站公告发布、内容审核；
		\item 系统运维监控模块：在线人数、判题队列、编译器负载、服务器健康监控、异常提交标记。
	\end{enumerate}
	
	\subsection{功能分配}
	\begin{itemize}
		\item 用户权限模块：承担登录校验、账号管理、接口鉴权全部功能；
		\item 题库管理模块：负责题目基础数据、测试用例、难度分类的存储与维护；
		\item 在线判题模块：承载代码接收、异步调度、沙箱执行、结果比对核心逻辑；
		\item 数据统计排行模块：读取提交记录，实时计算通过率、解题数量、排名数据；
		\item 讨论与公告模块：管理帖子、评论、公告内容及教师审核操作；
		\item 系统运维监控模块：采集服务器与业务运行指标，提供可视化监控视图与违规标记功能。
	\end{itemize}
	
	\section{接口设计}
	
	\subsection{接口架构与设计原则}
	本系统采用B/S架构和前后端分离模式。面向实训项目规模，后端优先采用\textbf{模块化单体}，将六个业务模块部署在同一Spring Boot应用中，通过清晰的应用服务接口隔离职责；文档生成属于CPU、内存和文件I/O密集型任务，采用独立工作进程异步执行。该方案兼顾小团队开发效率、事务一致性与后续拆分能力，不在当前阶段引入微服务治理成本。
	
	技术基线如下：
	\begin{itemize}
		\item 展示层：Vue 3、TypeScript、Vite和Element Plus构建单页应用；富文本编辑器采用基于ProseMirror内核的成熟组件，统一维护章节、段落、图片、表格和引用锚点；
		\item 接入层：Nginx提供HTTPS终止、静态资源发布、反向代理、下载限速和统一访问日志；
		\item 应用层：Spring Boot提供RESTful API，Spring Security完成身份认证与角色授权，业务代码按认证、学院、模板、写作、预览导出、批阅六个模块组织；
		\item 数据层：MySQL 8.0保存强一致业务数据，Redis保存短期会话状态、自动保存防抖键、编辑锁和异步任务状态；
		\item 文件层：对象存储保存论文图片、生成中的临时文件和可下载文件；开发环境可采用受控本地目录，生产环境建议使用MinIO等私有对象存储；
		\item 文档生成层：以“模板版本快照+论文结构化内容+媒体资源+参考文献”为统一输入，使用docx4j或Apache POI生成DOCX，再由LibreOffice Headless转换为PDF；网页预览、DOCX和PDF共享同一份规范化文档模型。
	\end{itemize}
	
	接口设计遵循以下原则：
	\begin{enumerate}
		\item \textbf{边界稳定}：跨模块仅传递业务标识、状态和必要的只读数据，不允许一个模块直接修改另一个模块拥有的数据表；
		\item \textbf{同步与异步分离}：登录、查询、保存、状态流转采用同步接口；预览重建、DOCX/PDF生成、过期文件清理采用异步任务；
		\item \textbf{版本可追溯}：论文始终绑定创建时的模板版本，提交、批阅和导出均携带论文内容版本，避免模板更新或并发保存改变历史结果；
		\item \textbf{统一鉴权}：所有业务接口在接入层之后仍由后端校验角色和资源归属，前端隐藏按钮不能替代权限控制；
		\item \textbf{失败可恢复}：会改变状态的接口具备幂等控制，异步任务可重试且保留失败原因，跨存储操作具备补偿和清理机制。
	\end{enumerate}
	
	\subsection{用户界面接口}
	系统仅面向PC端浏览器提供Web界面，不提供专用硬件接口和移动端原生接口。界面层以角色工作台为入口，并通过统一的导航、表单校验、确认弹窗和Toast反馈保持交互一致。
	
	\subsubsection{管理员界面}
	管理员端包含学院管理、模板列表、模板新增/编辑、模板详情和版本信息等页面。模板编辑页面按基本信息、封面字段、章节结构、排版参数四个区域组织。删除学院、删除模板、停用模板等操作必须显示影响范围并进行二次确认；存在被引用数据时，界面应阻止物理删除并提示采用停用方式。
	
	\subsubsection{学生界面}
	学生端包含论文工作台、模板选择、分章节编辑、参考文献管理、全文预览、导出任务和批阅结果等页面。编辑器持续显示“已保存、保存中、保存失败、离线”四类状态；自动保存失败时不得清空本地编辑内容。论文提交后，编辑控件切换为只读状态，并清晰展示当前状态、提交教师和最近一次批阅结果。
	
	\subsubsection{教师界面}
	教师端包含待批阅列表、论文只读预览、批注面板、整体评语、评分和历史记录。教师批注必须绑定到某次不可变的提交版本；当学生重新提交时，历史批注保持可查，新旧版本不得混写。
	
	\subsection{外部软件接口}
	系统无教务、查重、AI写作等第三方业务系统接口。对外软件接口主要包括浏览器访问接口、文件上传下载接口以及基础设施适配接口。
	
	\begin{table}[htbp]
		\centering
		\small
		\caption{主要外部接口分组}
		\begin{tabular}{p{2.3cm}p{3.0cm}p{3.0cm}p{4.2cm}}
			\toprule
			接口分组 & 调用方/提供方 & 主要职责 & 关键约束 \\
			\midrule
			认证与账户 & Web前端/应用服务 & 注册、登录、刷新登录状态、退出 & 密码不回传；限流；返回角色与最小必要资料 \\
			学院与模板 & 管理员端/应用服务 & 学院维护、模板筛选、配置、版本与上下架 & 仅管理员；模板保存生成新版本；引用中的版本不可破坏 \\
			论文写作 & 学生端/应用服务 & 创建论文、加载章节、保存内容、管理引用 & 仅论文所有者；内容版本校验；提交状态禁止写入 \\
			媒体与文献 & 学生端/文件及应用服务 & 图片上传、媒体引用、文献录入和排序 & 文件类型及大小校验；资源归属校验；引用序号统一重算 \\
			预览与导出 & 学生端/文档生成服务 & 创建预览、发起导出、查询进度、下载结果 & 基于已保存版本；任务幂等；下载地址短时有效 \\
			提交与批阅 & 学生端、教师端/应用服务 & 提交锁定、批注、评分、退回、通过、查看历史 & 仅课程论文；教师归属校验；状态机与事务控制 \\
			\bottomrule
		\end{tabular}
	\end{table}
	
	文件上传接口使用\texttt{multipart/form-data}，支持JPG、PNG、GIF和经过安全处理的SVG，单文件不超过5MB。服务端同时校验扩展名、MIME类型和文件特征，不以客户端声明作为唯一依据。下载接口不暴露对象存储物理路径，而是由应用服务校验权限后返回短时下载凭证或通过Nginx内部转发。
	
	\subsection{内部模块接口}
	六个业务模块及平台支撑组件之间的协作关系如下。
	
	\begin{table}[htbp]
		\centering
		\small
		\caption{内部模块协作关系}
		\begin{tabular}{p{2.5cm}p{3.0cm}p{3.0cm}p{3.8cm}}
			\toprule
			调用模块 & 被调用模块 & 交互内容 & 约束 \\
			\midrule
			全部业务模块 & 认证与权限模块 & 当前用户、角色、资源访问判定 & 每次敏感操作重新鉴权，不信任前端角色参数 \\
			模板管理模块 & 学院管理模块 & 校验学院存在与可用状态 & 模板必须归属有效学院，删除学院前检查引用 \\
			在线写作模块 & 模板管理模块 & 获取启用模板及指定版本快照 & 创建论文后固定版本，不跟随模板后续修改 \\
			在线写作模块 & 预览导出模块 & 提交论文版本、章节、资源和导出范围 & 只读取已保存内容，不直接修改论文 \\
			批阅模块 & 在线写作模块 & 读取提交快照并驱动论文状态变更 & 状态变更在事务中完成，批注绑定提交版本 \\
			预览导出模块 & 文件存储适配器 & 读取图片并写入生成文件 & 使用逻辑资源标识，不在业务层拼接物理路径 \\
			全部业务模块 & 审计与监控组件 & 记录操作事件、异常和性能指标 & 审计写入失败不得泄露业务数据，关键操作需告警 \\
			\bottomrule
		\end{tabular}
	\end{table}
	
	\subsection{异步任务接口}
	文档生成任务通过Redis任务队列或可靠任务表投递，任务内容仅保存任务标识、论文标识、论文内容版本、模板版本、导出格式、导出范围和发起人标识，不在消息体中复制整篇论文。工作进程领取任务后从MySQL和对象存储读取一致性快照，生成文件后回写任务状态与文件标识。
	
	导出任务状态统一为“等待、处理中、成功、失败、已过期”。同一用户对同一论文版本、同一格式和同一范围的重复请求复用未过期结果；工作进程异常退出后，超时的处理中任务可被重新领取。任务至少投递一次，因此生成过程必须以任务标识保证幂等。
	
	\subsection{通用接口约定}
	\begin{enumerate}
		\item 接口根路径统一使用版本前缀\texttt{/api/v1}，传输协议为HTTPS，字符编码为UTF-8，普通请求和响应使用JSON；
		\item 登录成功后使用短时访问令牌访问接口，刷新令牌仅用于换取新访问令牌；退出登录或账号禁用后，相关令牌进入失效状态；
		\item 响应统一包含业务结果、可读提示、追踪标识\texttt{traceId}和服务器时间；失败响应不返回堆栈、SQL或物理文件路径；
		\item 列表接口统一支持页码、每页数量、排序和白名单筛选条件，每页数量设置上限，防止一次读取过多数据；
		\item 保存接口携带内容版本号，服务端采用乐观锁判断并发冲突；状态变更接口携带幂等键，重复调用返回同一业务结果；
		\item HTTP状态码表达传输层结果，业务错误码表达认证、参数、权限、状态冲突、存储、生成和系统异常等类别。
	\end{enumerate}
	
	\section{数据结构设计}
	
	\subsection{数据架构原则}
	数据结构围绕“模板版本化、论文分章节、提交快照不可变、批阅可追溯、文件与业务记录分离”设计。MySQL作为事实数据源，Redis只保存可重建的短期数据，对象存储只保存二进制资源和生成文件。任何缓存或文件缺失均不得导致论文核心文本、模板版本和批阅结论丢失。
	
	\subsection{逻辑结构设计}
	\subsubsection{组织与身份数据}
	\begin{itemize}
		\item \textbf{用户}：保存用户名、密码摘要、角色、所属学院、账号状态和审计时间。用户名唯一；角色限定为管理员、学生、教师；
		\item \textbf{学院}：保存学院名称、描述和有效状态，为模板分类和用户归属提供基础数据；
		\item \textbf{教师分配关系}：保存课程论文与授课教师的对应关系，用于教师资源权限判定，不允许仅凭前端传入教师标识建立授权。
	\end{itemize}
	
	\subsubsection{模板数据}
	\begin{itemize}
		\item \textbf{论文模板}：保存模板名称、类型、学院、启停状态、当前版本号和说明，表示可管理的模板主记录；
		\item \textbf{模板版本}：保存某一版本的完整快照，包括封面字段、章节结构、排版参数和固定文本。配置采用带模式版本号的JSON存储，保存后不可原位覆盖；
		\item \textbf{模板资源}：保存模板引用的校徽、声明页背景等资源标识，不在JSON中保存二进制内容。
	\end{itemize}
	
	\subsubsection{写作与引用数据}
	\begin{itemize}
		\item \textbf{论文}：保存论文标题、学生、论文类型、绑定的模板版本、当前状态、内容版本、教师和时间信息；
		\item \textbf{论文章节}：按模板章节顺序保存章节类型、标题、层级、可编辑属性和富文本内容。内容以净化后的HTML或结构化JSON保存，并为段落、表格、图片等块分配稳定锚点；
		\item \textbf{参考文献}：保存文献类型、作者、标题、来源、年份、卷期和页码等结构化字段，并保存论文内排序号；
		\item \textbf{正文引用}：关联论文、章节锚点和参考文献，导出时根据文献顺序统一生成编号；
		\item \textbf{媒体资源}：保存原文件名、对象存储键、内容类型、大小、摘要值、上传者和引用状态，二进制数据不进入MySQL。
	\end{itemize}
	
	\subsubsection{导出与批阅数据}
	\begin{itemize}
		\item \textbf{导出任务}：保存论文内容版本、模板版本、格式、范围、任务状态、失败原因、生成文件和过期时间；
		\item \textbf{提交版本}：学生提交时固化论文标题、章节内容、引用关系、资源清单和模板版本，形成不可变快照；
		\item \textbf{批阅任务}：关联提交版本、学生和教师，记录待批阅、已通过、已退回等状态；
		\item \textbf{批注}：关联提交版本和稳定内容锚点，保存批注文本、创建人和创建时间；
		\item \textbf{评阅结果}：保存整体评语、分数、等级、结论和处理时间；等级由分数区间校验；
		\item \textbf{操作审计}：记录主体、动作、目标类型、目标标识、结果、IP、追踪标识和时间，不保存密码、令牌和论文全文。
	\end{itemize}
	
	\subsection{核心实体关系}
	\begin{enumerate}
		\item 一个学院可关联多个模板和多个用户；学院存在有效模板时不得直接物理删除；
		\item 一个模板包含多个不可变模板版本，模板主记录指向当前版本；一篇论文只绑定一个模板版本；
		\item 一个学生可创建多篇论文，一篇论文包含多个有序章节、参考文献、引用和媒体资源；
		\item 一篇课程论文可以产生多个提交版本，每次提交版本对应一次批阅任务；批注和评阅结果均绑定具体提交版本；
		\item 一个提交版本只能来源于一个已保存论文内容版本，历史提交内容不随学生后续修改而变化；
		\item 一个导出任务对应一个确定的论文内容版本和模板版本，生成文件可在有效期内被多次下载。
	\end{enumerate}
	
	\subsection{状态模型与一致性约束}
	\subsubsection{模板状态}
	模板状态为“已启用”或“已停用”。只有已启用模板可用于新建论文；停用、删除模板不影响已创建论文。模板编辑采用新增版本方式，旧版本不可覆盖，并通过模板版本号和配置模式版本号共同保证可解释性。
	
	\subsubsection{论文状态}
	论文状态按以下路径流转：
	\begin{itemize}
		\item 正常完成路径：草稿 $\rightarrow$ 已提交 $\rightarrow$ 已批阅；
		\item 退回修改路径：草稿 $\rightarrow$ 已提交 $\rightarrow$ 已退回 $\rightarrow$ 已提交。
	\end{itemize}
	只有草稿和已退回状态允许学生编辑；提交操作在同一事务内完成内容版本校验、提交快照创建、论文锁定和批阅任务创建。教师通过或退回时必须校验当前批阅任务仍处于待处理状态，防止重复处理。
	
	\subsubsection{并发与版本约束}
	论文和模板主记录设置乐观锁版本。学生保存时携带上次读取的内容版本，版本不一致则返回冲突并保留客户端草稿，由用户选择刷新或复制未保存内容。自动保存与手动保存使用相同的保存通道，按论文标识串行化短时间内的写请求，避免乱序覆盖。
	
	\subsection{物理结构设计}
	\begin{table}[htbp]
		\centering
		\small
		\caption{数据存储分工}
		\begin{tabular}{p{2.4cm}p{4.2cm}p{5.4cm}}
			\toprule
			存储组件 & 数据范围 & 设计要求 \\
			\midrule
			MySQL 8.0 & 用户、学院、模板及版本、论文及章节、文献、提交与批阅、任务与审计 & InnoDB事务；外键或等价完整性校验；业务唯一索引；时间字段统一；关键查询建立组合索引 \\
			Redis & 登录失效标识、自动保存防抖键、编辑短锁、导出任务进度和热点模板摘要 & 设置过期时间；不保存唯一副本；键包含环境和业务前缀；故障时允许降级到数据库 \\
			对象存储 & 图片、模板资源、DOCX/PDF成品和临时中间文件 & 私有桶；随机对象键；元数据入库；临时文件和过期导出按策略清理 \\
			\bottomrule
		\end{tabular}
	\end{table}
	
	关系表使用统一主键、创建时间、更新时间和必要的逻辑删除标识。模板版本、提交版本、评阅结果和审计记录属于历史证据数据，不执行普通业务删除。大文本章节与高频列表字段分离，列表页不加载论文全文。JSON配置必须通过服务端模式校验后写入，禁止将无法解释的任意字段直接落库。
	
	\subsection{索引与查询策略}
	\begin{itemize}
		\item 用户名、学院名称建立唯一索引；模板按学院、类型、启用状态和更新时间建立组合索引；
		\item 论文按学生、状态和更新时间建立组合索引；待批阅任务按教师、状态和提交时间建立组合索引；
		\item 章节按论文和顺序号唯一定位；参考文献按论文和排序号唯一定位；
		\item 导出任务按发起人、论文、内容版本、格式和状态查询；审计日志按时间、主体和目标标识检索；
		\item 论文预览所需章节、文献和资源采用批量读取，避免逐章节、逐资源产生大量数据库往返。
	\end{itemize}
	
	\subsection{数据结构与模块关系}
	\begin{itemize}
		\item 认证与权限模块拥有用户数据，对其他模块只提供身份和授权判定；
		\item 学院管理模块拥有学院数据，模板管理模块通过学院查询接口引用；
		\item 模板管理模块拥有模板、模板版本和模板资源，在线写作模块只读取指定版本快照；
		\item 在线写作模块拥有论文、章节、文献、引用和媒体元数据，预览导出模块只读取一致性快照；
		\item 批阅模块拥有提交版本、批阅任务、批注和评阅结果，并通过论文状态接口驱动锁定或解锁；
		\item 平台支撑组件拥有导出任务、审计和监控数据，不反向修改业务内容。
	\end{itemize}
	
	\section{运行设计}
	
	\subsection{部署与进程结构}
	系统在逻辑上分层，在物理部署上由五类运行单元组成：
	\begin{enumerate}
		\item \textbf{Web前端与Nginx}：发布静态资源、代理API、限制上传大小、转发受控下载；
		\item \textbf{应用服务}：运行六个业务模块，承担认证、事务、权限、状态流转和任务投递；
		\item \textbf{文档生成工作进程}：消费导出任务，装配规范化文档模型，生成DOCX并转换PDF；
		\item \textbf{数据服务}：MySQL提供持久化，Redis提供短期状态与队列；
		\item \textbf{文件服务}：对象存储保存上传资源和生成文件。
	\end{enumerate}
	
	开发及小规模验收环境可部署为单机多进程；生产环境建议将应用服务、文档生成进程、数据库和对象存储分离。应用服务应保持无状态，可部署多个实例；生成进程按CPU和内存容量独立扩缩，不与在线保存接口争用同一工作线程池。
	
	\subsection{运行模块组合}
	\begin{table}[htbp]
		\centering
		\small
		\caption{主要业务场景的模块组合}
		\begin{tabular}{p{2.4cm}p{5.0cm}p{4.6cm}}
			\toprule
			业务场景 & 参与模块 & 协作结果 \\
			\midrule
			模板发布 & 认证权限、学院管理、模板管理、审计组件 & 校验学院和管理员权限，生成模板新版本并切换启用状态 \\
			创建与写作 & 认证权限、模板管理、在线写作、文件存储 & 绑定模板版本，创建章节，保存正文、文献和图片 \\
			预览与导出 & 在线写作、预览导出、文档生成进程、对象存储 & 读取已保存快照，生成预览或DOCX/PDF并提供受控下载 \\
			提交与批阅 & 认证权限、在线写作、批阅、审计组件 & 固化提交版本并锁定，教师批注评分，按结论解锁或完成 \\
			运维恢复 & 监控审计、数据库、Redis、对象存储 & 定位失败任务、恢复服务、重放幂等任务并验证数据完整性 \\
			\bottomrule
		\end{tabular}
	\end{table}
	
	\subsection{关键运行流程}
	\subsubsection{论文自动保存流程}
	学生编辑产生变更后，前端先更新本地草稿，并在空闲窗口内合并变更；每30秒或用户手动保存时，将论文标识、内容版本和变化后的章节内容提交后端。后端依次完成身份校验、资源归属校验、论文状态校验、内容净化、字数与资源限制校验、乐观锁更新和保存审计。保存成功返回新内容版本；网络失败时前端保留本地草稿并提示重试。
	
	\subsubsection{预览与导出流程}
	用户发起预览或导出时，系统只使用最近一次保存的内容版本。应用服务建立导出任务并立即返回任务标识，文档生成进程异步执行以下步骤：
	\begin{enumerate}
		\item 读取论文、模板版本、章节、参考文献和资源清单，验证快照完整性；
		\item 将模板配置和论文内容合并为规范化文档模型，统一计算章节编号、目录、引用编号和分页属性；
		\item 预览适配器生成分页HTML；Word适配器生成DOCX；PDF适配器由同一DOCX转换生成PDF；
		\item 校验生成文件存在、大小非零、页数可读取后写入对象存储；
		\item 更新任务状态并记录生成耗时、文件摘要和失败原因。
	\end{enumerate}
	前端通过任务查询接口获取进度，不保持长连接占用应用线程。任务失败允许用户重试，但已成功且未过期的相同任务直接复用。
	
	\subsubsection{提交与批阅流程}
	学生确认提交后，后端在事务内校验论文类型、教师归属、必填章节和最新内容版本，创建不可变提交快照，将论文状态改为已提交并生成待批阅任务。教师打开的始终是该提交快照，不是学生工作副本。教师保存批注可多次进行，最终通过或退回时一次性校验分数、等级和状态；通过后论文进入已批阅，退回后进入已退回并允许学生继续编辑。
	
	\subsection{运行控制}
	\begin{enumerate}
		\item \textbf{认证控制}：Nginx只负责接入，应用服务对每个接口执行令牌、角色和资源归属校验；
		\item \textbf{事务控制}：模板版本发布、论文提交、教师结论等多记录变更在单个数据库事务内完成；
		\item \textbf{并发控制}：模板和论文使用乐观锁；导出任务使用任务租约；同一批阅任务只允许一个终态写入；
		\item \textbf{幂等控制}：提交、通过、退回和导出请求携带幂等键，重复请求不得产生多份业务记录；
		\item \textbf{超时控制}：数据库、Redis、对象存储和文档转换均设置连接及执行超时，超时后释放资源并记录追踪标识；
		\item \textbf{限流控制}：登录、图片上传和导出接口按用户及IP限流；导出任务按用户设置并发上限；
		\item \textbf{降级控制}：Redis不可用时登录和数据库业务可继续运行，自动保存防抖和热点缓存退化为数据库方案；文档服务不可用时编辑与保存仍可使用。
	\end{enumerate}
	
	\subsection{性能与容量设计}
	性能目标与需求规格保持一致：
	\begin{itemize}
		\item 标准网络环境下普通页面首屏加载不超过2秒，普通查询接口的服务端P95响应时间控制在1秒以内；
		\item 论文保存操作服务端处理时间不超过1秒，自动保存不得阻塞编辑器输入；
		\item 网页预览在3秒内可用；万字以内论文的DOCX或PDF导出目标不超过10秒，超出时仍以异步任务方式完成并展示状态；
		\item 单篇论文字数不超过50000字，图片不超过50张且单张不超过5MB，参考文献不超过200条；
		\item 应用服务线程池与文档生成线程池隔离。文档生成并发数依据单任务峰值内存和CPU核数配置，不以无限排队换取表面成功率；
		\item 列表接口默认分页，论文正文按需加载；模板摘要可缓存，论文正文、提交快照和评阅结果不依赖缓存作为唯一数据源。
	\end{itemize}
	
	\subsection{定时任务}
	\begin{itemize}
		\item 每5分钟扫描超时的导出任务，将失去租约的任务恢复为可重试状态；
		\item 每日清理超过有效期且无业务引用的导出文件、上传临时文件和失败中间文件；
		\item 每日校验数据库记录与对象存储对象的引用一致性，输出孤儿文件和缺失文件清单；
		\item 每日执行数据库备份并校验备份结果，定期进行恢复演练；
		\item 每月归档过期访问日志与运行指标，审计日志按安全策略保留。
	\end{itemize}
	
	\section{出错处理设计}
	
	\subsection{处理原则与异常分类}
	系统按照“对用户可理解、对运维可定位、对业务可恢复”的原则处理异常。异常分为参数校验、认证授权、资源不存在、业务状态冲突、并发冲突、文件处理、基础设施、文档生成和未知系统异常九类。每次失败均生成追踪标识，用户提示与后台日志通过该标识关联。
	
	\subsection{出错输出信息}
	\begin{enumerate}
		\item 用户侧只显示可执行的提示，例如“内容版本已变化，请刷新后重试”“图片超过5MB”“论文已提交，当前不可编辑”“导出失败，可重新发起”；
		\item 接口响应包含稳定错误码、简短消息、必要的字段级校验信息和\texttt{traceId}，不包含堆栈、SQL、服务器目录、对象存储键和安全配置；
		\item 后台日志记录异常类型、业务标识、调用链、耗时和受控上下文。论文正文、密码、令牌和上传文件内容不得写入日志；
		\item 异步任务失败记录失败阶段、可重试性、重试次数和最后异常摘要，前端可展示任务失败状态但不展示内部异常详情。
	\end{enumerate}
	
	\subsection{主要异常及处理对策}
	\begin{table}[htbp]
		\centering
		\small
		\caption{异常处理策略}
		\begin{tabular}{p{2.5cm}p{3.5cm}p{3.6cm}p{2.8cm}}
			\toprule
			异常场景 & 用户侧处理 & 服务端处理 & 恢复与告警 \\
			\midrule
			登录失败或越权 & 提示凭据错误、账号禁用或权限不足 & 统一拒绝并记录安全事件；不区分用户名是否存在 & 短时间多次失败触发限流和告警 \\
			保存版本冲突 & 保留本地草稿，提示刷新或复制内容后处理 & 拒绝覆盖，返回当前版本摘要 & 统计冲突率，检查多端同时编辑情况 \\
			网络中断 & 显示离线状态，保留本地编辑内容 & 未收到完整请求不写库；幂等请求可重试 & 网络恢复后由用户或客户端重试 \\
			图片上传失败 & 明确提示格式、大小或网络原因 & 删除未完成临时对象，保持论文内容不变 & 重试失败或安全检测异常时告警 \\
			数据库或Redis异常 & 提示暂时不可用，编辑内容留在本地 & 数据库事务回滚；Redis功能降级，不写入不一致结果 & 熔断、健康检查、连接恢复后验证 \\
			文档生成失败 & 保留论文，显示失败任务并允许重试 & 标记失败阶段，清理中间文件，不覆盖旧成功文件 & 达到重试上限后人工介入并告警 \\
			对象存储不可用 & 暂停上传和下载，写作文本仍可保存 & 快速失败，不生成无文件的成功记录 & 恢复后重试任务并执行引用一致性检查 \\
			\bottomrule
		\end{tabular}
	\end{table}
	
	\subsection{重试、补偿与回滚}
	\begin{enumerate}
		\item 只对查询、幂等写入和明确标记可重试的异步任务自动重试，采用递增退避并设置最大次数；认证失败、参数错误和业务状态冲突不自动重试；
		\item 数据库事务失败全部回滚。数据库记录与对象存储无法组成分布式事务时，先上传到临时区，再提交业务记录，最后转为正式对象；任一步失败均由清理任务补偿；
		\item 导出任务生成新文件成功并校验后才更新文件引用，失败时保留上一份成功文件，不向用户提供半成品；
		\item 论文提交失败不得产生“已锁定但无提交版本”或“有批阅任务但仍可编辑”的中间状态；启动时和定时任务均检查此类不变量；
		\item 教师最终结论写入失败时，批注草稿可以保留，但论文状态和评阅终态必须保持未变，允许教师重新提交结论。
	\end{enumerate}
	
	\subsection{日志、监控与告警}
	应用日志采用结构化格式，并携带请求追踪标识、用户匿名标识、模块、接口、耗时和结果。监控至少覆盖接口错误率与P95耗时、数据库连接池、Redis可用性、导出队列长度、最老任务等待时间、文档生成成功率、对象存储容量和磁盘临时目录使用率。
	
	当普通接口连续5分钟错误率超过1\%、保存接口P95超过1秒、最老导出任务等待超过2分钟、文档生成连续失败、存储使用率超过80\%或数据库连接池持续超过80\%时触发告警。告警应包含环境、时间、模块、指标和值，不包含论文正文和个人敏感信息。
	
	\section{安全保密设计}
	
	\subsection{安全目标与信任边界}
	系统保护的核心资产包括账号凭据、用户个人信息、论文内容、模板配置、提交快照、教师批注与评分、上传图片和导出文件。浏览器、互联网请求和用户上传内容均视为不可信；Nginx之后的应用服务仍需进行完整校验；数据库、Redis和对象存储仅允许内网访问。
	
	\subsection{身份认证与权限控制}
	\begin{enumerate}
		\item 密码采用BCrypt或Argon2等带盐单向摘要算法保存，禁止明文、可逆加密和日志记录；
		\item 登录采用短时访问令牌和受控刷新机制，令牌包含最小必要身份信息；账号禁用、退出和密码变更后使相关令牌失效；
		\item 采用RBAC与资源归属双重授权：管理员仅管理学院和模板；学生仅管理本人论文；教师仅查看并批阅分配给本人的课程论文；
		\item 后端根据当前登录主体确定用户标识，不接受客户端通过修改学生、教师或管理员标识绕过授权；
		\item 论文提交、模板删除、模板停用、评分结论等敏感操作要求二次确认，并记录审计事件。
	\end{enumerate}
	
	\subsection{接口与会话安全}
	\begin{itemize}
		\item 对外通信强制HTTPS，建议TLS 1.2及以上；HTTP自动重定向到HTTPS；
		\item CORS仅允许受控前端域名；根据令牌存放方式配置CSRF防护，并设置CSP、X-Content-Type-Options等安全响应头；
		\item 登录、注册、上传和导出接口实施用户及IP维度限流；异常请求不消耗无限线程、内存或临时磁盘；
		\item 所有查询采用参数化访问，排序字段和筛选字段使用白名单，防止SQL注入和任意字段访问；
		\item 服务端统一校验长度、枚举、数字范围、状态和业务归属，不能仅依赖前端表单校验。
	\end{itemize}
	
	\subsection{富文本与文件安全}
	\begin{enumerate}
		\item 富文本内容使用标签、属性和协议白名单净化，移除脚本、事件属性、外部执行内容和危险URL，预防存储型XSS；
		\item 图片上传同时校验扩展名、MIME类型、文件特征、大小和数量；文件名只作为展示信息，存储键由系统随机生成；
		\item SVG在入库前删除脚本、外链、事件属性和嵌入对象，无法安全净化时转换为PNG或拒绝上传；
		\item 上传对象放入私有存储区域，禁止直接执行；下载时设置正确的内容类型和Content-Disposition；
		\item DOCX/PDF生成在低权限工作目录中完成，限制单任务CPU、内存、执行时间和临时磁盘，任务结束后清理临时文件；
		\item 文档生成引擎不得加载用户提供的宏、外部链接资源或任意本地路径，模板中所有资源必须来自已登记的对象标识。
	\end{enumerate}
	
	\subsection{数据保护与保密}
	\begin{itemize}
		\item 数据库账号遵循最小权限，应用、备份和运维使用不同账号；Redis和对象存储启用认证并限制网络来源；
		\item 生产密钥、数据库密码和对象存储凭据通过环境密钥或专用配置中心注入，不写入源代码和版本库；
		\item 备份文件、导出文件和传输链路按环境能力加密；下载文件使用短时授权，过期后自动失效；
		\item 页面和接口仅返回完成当前业务所需的字段；教师不可读取未提交论文，管理员不可读取学生正文；
		\item 论文批阅历史属于业务证据数据，按论文生命周期保留；临时导出文件按有效期清理，用户注销或数据清理按学校制度执行。
	\end{itemize}
	
	\subsection{审计与安全事件}
	模板新增、版本发布、启停和删除，论文提交、退回、通过和评分，账号禁用，越权访问，连续登录失败，文件安全校验失败等操作必须进入审计日志。审计记录应防止普通业务用户修改，包含操作人、角色、目标、结果、时间、来源IP和追踪标识。
	
	出现疑似泄露或恶意访问时，运维人员可按追踪标识冻结相关账号和令牌、暂停导出、保存证据日志、核查访问范围并恢复服务。安全事件处置过程不得通过删除业务数据掩盖问题。
	
	\section{维护设计}
	
	\subsection{维护目标与总体策略}
	维护设计以“可定位、可替换、可回滚、可恢复”为目标。当前系统采用模块化单体而非微服务，六个业务模块拥有明确数据和服务边界；文档生成、文件存储、缓存等变化较快或资源特征明显的能力通过适配器隔离，使小团队能够在不改变核心业务流程的前提下升级技术组件。
	
	\subsection{模块与依赖维护}
	\begin{enumerate}
		\item 业务模块只依赖公开的应用服务接口，不跨模块直接访问数据访问对象；公共代码限于认证上下文、错误模型、审计、时间和基础工具；
		\item 文档生成采用“规范化文档模型+格式适配器”，新增导出格式或替换DOCX/PDF引擎时不修改写作和批阅模块；
		\item 文件存储通过统一资源接口访问，开发环境本地存储和生产对象存储可按配置切换；
		\item 模板JSON包含模式版本号，新增排版参数时提供向后兼容默认值和升级程序，历史模板版本必须仍可导出；
		\item 数据库变更使用Flyway或同类迁移工具按版本执行，迁移脚本只追加不篡改，发布前在备份副本上验证。
	\end{enumerate}
	
	\subsection{配置与版本管理}
	开发、测试和生产环境分别维护配置，数据库地址、Redis地址、对象存储、令牌有效期、上传限制、导出并发、文件有效期和告警阈值全部外置。敏感配置与普通配置分离，生产密钥不得出现在配置样例中。
	
	前端、应用服务、文档生成进程和数据库模式均记录版本。导出文件元数据保存生成引擎版本、模板版本和论文内容版本，出现排版问题时可复现生成条件。API采用主版本前缀；不兼容变更发布新版本并保留迁移窗口。
	
	\subsection{日志与监控维护}
	\begin{itemize}
		\item 访问日志、应用日志、文档生成日志、审计日志分流保存，统一时间、环境、级别、模块和追踪标识；
		\item 访问及应用日志按天滚动，在线保留30天后归档或清理；审计日志至少保留180天，具体期限服从学校制度；
		\item 监控面板展示服务可用性、接口耗时、保存成功率、导出任务、数据库、Redis、对象存储和主机资源；
		\item 每周复核慢查询、重复异常和失败任务；对高频问题建立故障知识库，记录现象、影响、定位方法和恢复步骤；
		\item 日志清理策略必须先于磁盘占满，临时目录、导出目录和数据库日志分别设置容量阈值。
	\end{itemize}
	
	\subsection{数据备份与恢复}
	\begin{enumerate}
		\item MySQL每日全量备份并持续保留二进制日志，目标恢复点RPO不超过15分钟；模板版本、提交版本、评阅结果和审计数据纳入核心备份；
		\item 对象存储每日增量备份论文图片和模板资源；可重新生成且已过期的导出文件不作为最高优先级备份对象；
		\item Redis不承担唯一数据存储，故障后允许从数据库重建缓存和任务状态；
		\item 备份文件异地或跨磁盘保存、加密并校验摘要；仅有备份文件不视为备份成功，至少每季度执行一次恢复演练；
		\item 目标恢复时间RTO不超过2小时。恢复顺序为数据库、对象存储、应用服务、文档生成进程、缓存，恢复后执行数量、引用和抽样导出校验。
	\end{enumerate}
	
	\subsection{发布、回滚与兼容性维护}
	发布流程依次执行代码检查、单元测试、接口测试、模板兼容测试、典型论文导出回归、数据库迁移验证和安全扫描。构建产物应不可变并带版本号。生产发布前完成备份，优先采用双实例滚动或双槽切换，确认健康检查、保存和导出均正常后再结束观察期。
	
	回滚时先回退应用和前端产物；数据库迁移优先采用向前兼容设计，避免依赖破坏性回滚。涉及文档引擎升级时保留上一版本工作进程和典型样例基线，一旦出现严重排版差异，可将新任务切回上一引擎。
	
	\subsection{例行维护计划}
	\begin{table}[htbp]
		\centering
		\small
		\caption{例行维护项目}
		\begin{tabular}{p{2.0cm}p{6.0cm}p{4.0cm}}
			\toprule
			周期 & 维护内容 & 通过标准 \\
			\midrule
			每日 & 检查备份、告警、导出失败任务、磁盘与对象存储容量 & 备份成功；无未解释严重告警；积压任务已处置 \\
			每周 & 分析慢查询、错误趋势、账号异常和孤儿文件清单 & 关键问题有负责人、原因和完成期限 \\
			每月 & 安装经验证的依赖安全补丁，复核权限、日志容量和告警阈值 & 无高危未处置漏洞；权限符合最小化原则 \\
			每季度 & 执行备份恢复、故障切换和典型DOCX/PDF回归演练 & 满足RPO/RTO；典型文档排版基线通过 \\
			每次发布 & 执行自动化测试、数据库迁移验证和回滚预案检查 & 发布清单签字确认，观察期无阻断问题 \\
			\bottomrule
		\end{tabular}
	\end{table}
	
	\subsection{扩展与演进设计}
	\begin{itemize}
		\item 新增论文类型时，优先扩展模板类型、章节配置和校验规则，不复制整套写作流程；
		\item 新增导出格式时，实现新的文档格式适配器，并复用规范化文档模型、资源解析和任务管理；
		\item 未来对接教务系统时，在组织与身份边界增加同步适配层，不让外部字段直接进入核心领域模型；
		\item 当导出量成为主要瓶颈时，可独立扩容文档生成工作进程；当模块确有独立发布和容量需求时，再基于现有接口边界拆分服务；
		\item 任何演进均应保持历史模板版本、提交版本和评阅记录可读取、可预览、可导出。
	\end{itemize}
	
\end{document}
